\section{Introduction}
\paragraph{}
A processor is considered functionally correct when it successfully executes the complete set of target ISA instructions, producing results consistent with the architectural specification. To verify RV32ICore, a set of C programs of increasing complexity are compiled using the xPack RISC-V GCC toolchain and executed in simulation. Correct execution is determined by inspecting the values held in memory or the register file at the end of simulation and comparing them against analytically computed expected results. 
\section{Simulation Environment:}
Most of the tests performed on RV32ICore will follow a structured simulation environment. The environment contains:
\subsection{Testbench Code:}
\begin{minted}[baselinestretch=1]{vhdl}
	Library ieee;
	use ieee.std_logic_1164.all;
	entity Hazard_Testbench is
	end entity;
	
	architecture arch of Hazard_Testbench is
	component P_Processor_Hazard is
	port(...
	);
	end component;
	component RAM_Inst is
	port(...
	);
	end component;
	component RAM is
	port(...
	);
	end component;
	Signal clk ,reset: std_logic;
	Signal instruction, ReadData: std_logic_vector(31 downto 0);
	Signal MemWrite: std_logic;
	Signal ByteEn, Mode: std_logic_vector(1 downto 0);
	Signal Pc: std_logic_vector(31 downto 0);
	Signal DataAdr, WriteData: std_logic_vector(31 downto 0);
	Signal clk_cnt, Inst_Change_Cnt: integer := 0;--number Of clocks counter
	begin
	UUT: P_Processor_Hazard port map(...
	);
	Data_RAM: RAM port map(...
	);
	Inst_RAM: RAM_Inst port map(...
	);
	clock: process Begin
	loop
	clk<= '1';
	wait for 10 ns;
	clk<= '0';
	wait for 10 ns;
	end loop;
	end process;
	
	Clock_counter: process(clk) 
	variable last: std_logic_vector(31 downto 0);
	Begin
	if(rising_edge(clk)) then
	clk_cnt<= clk_cnt + 1;
	if (instruction /= last) then
	Inst_Change_Cnt <= Inst_Change_Cnt + 1;
	end if;
	last := instruction;
	end if;
	end process;
	
	Stimulus: process Begin
	reset<= '1';
	wait for 40 ns;
	reset<= '0';
	wait;
	end process;
	end architecture;
\end{minted}
\label{RV32ICore Testbench Code}
The code maps between the processor and its memories and defines the necessary signals to do so. The testbench initializes by applying a reset for two clock cycles ensuring that the Pc, register file, and memory are all starting with Zero. A clock of 20 ns period is simulated using a Clock process and wait commands, and a counter is added to calculate the number of instructions preformed in a program used to calculate performance.
\subsection{Memory Loading}
The programs used for testing can be very long making impractical to load them manually. A more efficient way is to use the compiler. The output file can be turned into an output hex file which can be directly loaded into memory using a command supported by ModelSim (mem load -infile output.hex -format hex \allowbreak  sim:/hazard\_testbench/Inst\_RAM/Mem\allowbreak).
\subsection{Running The Simulation}
By running the simulation for an appropriate amount of time the results generated by the processor can be verified in memory where they are stored. Additionally if a it is necessary to check the behavior of any signal during the processing it is possible to do so using Waveforms.
\section{Tests}
Two tests will be preformed one of a summing function called by main and the other is more complex containing multiple arithmetic functions.
\subsection{Test N:°1}

The Processor initialize by resetting for two clock periods the PC starts with address 0 and increment by 4 in Normal sequence (No execution flow instructions) instructions are being fetched as the waveform suggests.

\pagebreak

\begin{figure}[h!]
	\centering
	\includegraphics[width=0.7\linewidth]{C:/Users/dell/Pictures/Screenshots/Waveform}
	\caption{R32ICore Initialization}
	\label{fig:waveform}
\end{figure}

When the execution of the program ends we can check the final values stored in stack due to instruction
\begin{minted}{riscv}
	E4:	sw  a0,-20(s0) ; [000003BC] = a0 where a0= 16H
\end{minted}
The following Figure indicates that indeed the program executed as expected by computing the value 16H and storing it in location address 0x000003BC, which when divided by the word size of 4 bytes corresponds to word index 0xEF in the data memory.
\begin{figure}[h!]
	\centering
	\includegraphics[width=0.7\linewidth]{"C:/Users/dell/Pictures/Screenshots/Sum array result"}
	\caption{Sum Array Result}
	\label{fig:sum-array-result}
\end{figure}
With this result it is appropriate to say that the processor is working properly another test will be done for further verification.
\subsection{Test N:°2}
The second test uses the following C code. It contains the previously used array summer it also contains a multiplayer, a pointer manipulator, and a bit-wise manipulator. The analyzed values should be 16H, 78H, 09H, and 08H respectively. 
\begin{minted}[baselinestretch=1]{c}
	// test.c
	
	// Tests loops and addition
	int sum_array(int arr[], int len) {
		int sum = 0;
		for (int i = 0; i < len; i++) {
			sum += arr[i];
		}
		return sum;
	}
	
	// Tests recursion and stack depth
	int factorial(int n) {
		int result = 1; // multiply by repeated addition
		for (int i = 2; i <= n; i++) {
			int temp = 0;
			for (int j = 0; j < i; j++) {
				temp += result;  
			}
			result = temp;
		}
		return result;
	}
	// Tests pointer manipulation
	int find_max(int arr[], int len) {
		int max = arr[0];
		for (int i = 1; i < len; i++) {
			if (arr[i] > max)
			max = arr[i];
		}
		return max;
	}
	
	// Tests bitwise operations
	unsigned int count_bits(unsigned int n) {
		unsigned int count = 0;
		while (n) {
			count += n & 1;
			n >>= 1;
		}
		return count;
	}
	
	int main() {
		int arr[] = {3, 7, 2, 9, 1};
		int s = sum_array(arr, 5);       // expected: 22
		int f = factorial(5);             // expected: 120
		int m = find_max(arr, 5);        // expected: 9
		unsigned int b = count_bits(255); // expected: 8
		return 0;
	}
\end{minted}
By checking the contents of the memory we find.
\begin{figure}[h!]
	\centering
	\includegraphics[width=0.7\linewidth]{"C:/Users/dell/Pictures/Screenshots/Second Test"}
	\caption{Second Test Results}
	\label{fig:second-test}
\end{figure}

The processor is preforming as intended and further testing was done on the unused instructions by the previous programs. 
\section{Conclusion}
\paragraph{}
This chapter has presented the verification methodology and simulation results for RV32ICore. A structured testing approach was adopted in which C programs of increasing complexity were compiled using the xPack RISC-V GCC toolchain and executed in the ModelSim simulation environment. 
\paragraph{}
Correct execution was determined by comparing values stored in the data memory against analytically derived expected results, with complete assembly listings provided in the Appendix for reference. Reset behavior was verified to be correct, with the PC initializing to zero and the pipeline filling in proper sequence following reset release. Two principal test programs were used: the first verified loop execution, load and store operations, and basic arithmetic through a sum-of-array computation producing the correct result of 22 (0x16); the second subjected the processor to a more demanding workload involving multiple concurrent function calls, factorial computation, pointer manipulation, and bitwise operations, yielding the correct results of 22, 120, 9, and 8 respectively. The remaining RV32I instructions were individually verified through targeted test programs. The combined results confirm that RV32ICore correctly executes the full RV32I Base Integer ISA under simulation.